\ifx\isStandalone\undefined
\documentclass[a4paper,10pt]{article}

% --- PACCHETTI FONDAMENTALI ---
\usepackage[utf8]{inputenc}
\usepackage[T1]{fontenc}
\usepackage[italian]{babel}
\usepackage[margin=2.5cm]{geometry}
\usepackage{lmodern}

% --- PACCHETTI PER GRAFICA E COLORI ---
\usepackage{tikz}
\usetikzlibrary{calc}
\usepackage{xcolor}

% --- PACCHETTI PER I BLOCCHI DI CODICE ---
\usepackage[most]{tcolorbox}
\usepackage{listings}

% --- DEFINIZIONE DEI COLORI PERSONALIZZATI ---
\definecolor{paletteIndigo}{HTML}{303F9F}   % Indaco scuro
\definecolor{paletteLavender}{HTML}{C5CAE9} % Lavanda chiaro
\definecolor{paletteLenurple}{HTML}{7B1FA2} % Viola acceso
\definecolor{textGray}{HTML}{757575}        % Grigio per il testo

% --- CONFIGURAZIONE LINGUAGGIO JAVASCRIPT PER LISTINGS ---
\lstdefinelanguage{JavaScript}{
  keywords={typeof, new, true, false, catch, function, return, null, catch, switch, var, if, in, while, do, else, case, break, const, let, class, extends, static, throw},
  keywordstyle=\color{paletteIndigo}\bfseries,
  ndkeywords={class, export, boolean, throw, implements, import, this},
  ndkeywordstyle=\color{darkgray}\bfseries,
  identifierstyle=\color{black},
  sensitive=false,
  comment=[l]{//},
  morecomment=[s]{/*}{*/},
  commentstyle=\color{green!60!black}\ttfamily,
  stringstyle=\color{red}\ttfamily,
  morestring=[b]',
  morestring=[b]",
  literate={à}{{\`a}}1
           {è}{{\`e}}1
           {é}{{\'e}}1
           {ì}{{\`i}}1
           {ò}{{\`o}}1
           {ù}{{\`u}}1
}

% --- DEFINIZIONE DELL'AMBIENTE JSBOX ---
\newtcblisting{jsbox}[1][]{
  enhanced,
  colback=gray!5,
  colframe=paletteIndigo,
  fonttitle=\bfseries,
  listing only,
  listing options={
    language=JavaScript,
    basicstyle=\ttfamily\footnotesize,
    breaklines=true,
    showstringspaces=false
  },
  #1 % Permette di passare il parametro title={}
}

\begin{document}
\fi

\newpage
\begin{titlepage}
\thispagestyle{empty}
\begin{tikzpicture}[remember picture, overlay]
    % Sfondo Laterale (Viola scuro per stile Bibbia)
    \fill[paletteIndigo] (current page.north west) rectangle ($(current page.south west) + (4cm,0)$);
    
    % Linea Accento (Lavanda)
    \draw[paletteLavender, line width=4pt] ($(current page.north west) + (4.2cm,-2cm)$) -- ($(current page.south west) + (4.2cm,2cm)$);
    
    \node[anchor=north, inner sep=0, text width=12cm, align=center] at ($(current page.north west) + (12.5cm,-5cm)$) {
        {\fontsize{18}{22}\selectfont \textcolor{textGray}{Università di Pisa}}\\[0.2cm]
        {\fontsize{16}{20}\selectfont \textcolor{textGray}{Corso di Laurea in Informatica}}\\[2cm]
        
        {\fontsize{50}{60}\selectfont \textbf{\textcolor{paletteIndigo}{CHEAT SHEET}}}\\[0.5cm]
        
        {\fontsize{30}{40}\selectfont \textbf{LA BIBBIA DI JAVASCRIPT}}\\[1.5cm]
        
        \textcolor{paletteLenurple}{\rule{10cm}{4pt}}\\[1cm]
        
        {\Large \textbf{Riepilogo Definitivo (Diviso per Lezioni)}}\\[0.2cm]
        {\large \textit{Sintassi, Metodi, Algoritmi e Snippet Salva-Vita}}\\[2cm]
        
        \textbf{\large Autore del riassunto:}\\
        {\large Joseph Zucchelli (IlBroccolo)}\\[2cm]
        
        \vfill
        {\large \textbf{Anno Accademico 2025 - 2026}}
    };
    
    % Badge CS
    \node[anchor=south east] at ($(current page.south east) + (-2cm, 2cm)$) {
        \tikz\draw[fill=paletteLavender, draw=none] (0,0) rectangle (2,2) node[midway, paletteIndigo, font=\bfseries\huge] {CS};
    };
\end{tikzpicture}
\end{titlepage}
\clearpage

\section{Cheat Sheet: La Bibbia di JavaScript (Per Lezioni)}
Riepilogo dettagliato strutturato in base all'ordine progressivo del corso.

%============================
% PART 1: Lezioni 01-10
%============================
\subsection{Lezioni 01-02: Fondamenti, Tipi e Operatori}
\begin{itemize}
    \item \textbf{Strict Mode}: \texttt{"use strict";} all'inizio del file previene errori silenziosi (es. variabili non dichiarate).
    \item \textbf{Tipi Primitivi}: \texttt{Number} (tutti i numeri, inclusi float e \texttt{NaN}), \texttt{String}, \texttt{Boolean}, \texttt{undefined} (variabile dichiarata ma non inizializzata), \texttt{null} (valore nullo intenzionale).
    \item \textbf{Operatori Logici}: \texttt{\&\&} (AND), \texttt{||} (OR), \texttt{!} (NOT). Valutazione "Short-circuit": \texttt{a || b} restituisce \texttt{a} se \textit{truthy}, altrimenti \texttt{b}.
    \item \textbf{Uguaglianza}: Usare SEMPRE \texttt{===} e \texttt{!==} per evitare conversioni di tipo implicite (es. \texttt{1 == "1"} è \texttt{true}, \texttt{1 === "1"} è \texttt{false}).
\end{itemize}
\begin{jsbox}[title={Snippet: Template Literals e Operatori}]
// String interpolation
console.log(`Esattamente ${2 + 2} gatti`); 

// Operatore Ternario (if inline)
let msg = eta >= 18 ? "Maggiorenne" : "Minorenne";
\end{jsbox}

\subsection{Lezione 03: Flusso di Controllo}
\begin{jsbox}[title={Snippet: If, Switch e Cicli}]
// If/Else
if (x > 0) { /*...*/ } else if (x < 0) { /*...*/ } else { /*...*/ }

// Switch
switch(voto) { 
    case 10: console.log("Max"); break; 
    default: console.log("Altro"); 
}

// Cicli: for, while, do..while
for(let i = 0; i < 5; i++) { if (i === 2) continue; if (i === 4) break; }
while(n > 0) { n--; } // Valuta condizione PRIMA
do { n++; } while(n < 5); // Valuta CONDIZIONE DOPO (esegue almeno 1 volta)
\end{jsbox}

\subsection{Lezione 04-05: Strutture Dati Base (Oggetti e Array)}
\textbf{Oggetti}: Insiemi di coppie chiave-valore. Le chiavi sono \texttt{String} (o convertite a stringa).
\textbf{Array}: Oggetti speciali ottimizzati per indici numerici con proprietà \texttt{length} e metodi iteratori.

\begin{jsbox}[title={Snippet: Interazione con Array e Oggetti}]
let obj = { nome: "Aldo", "età": 30 };
console.log(obj.nome); // Dot notation
console.log(obj["età"]); // Bracket notation (obbligatoria se chiave con spazi/simboli o variabile)

// Array Mutanti (Attenzione!)
let a = [1, 2, 3];
a.push(4);    // Aggiunge in fondo -> [1,2,3,4]
a.pop();      // Rimuove e restituisce l'ultimo [4] -> [1,2,3]
a.unshift(0); // Aggiunge in testa -> [0,1,2,3]
a.shift();    // Rimuove e restituisce la testa [0] -> [1,2,3]

// Array Splice vs Slice
a.splice(1, 1); // Rimuove 1 elem da indice 1. a diviene [1, 3] (MUTANTE)
let b = a.slice(0, 1); // Ritorna copia da 0 escluso 1. b=[1], a è intatto (NON MUTANTE)
\end{jsbox}

\subsection{Lezioni 06-07: Funzioni, Scope e Closures}
Ogni funzione è un oggetto di primo livello.
\begin{itemize}
    \item \textbf{Hoisting}: \texttt{function f()\{\}} può essere chiamata prima della sua definizione testuale. \texttt{let f = function()\{\}} non gode di hoisting.
    \item \textbf{Scope (var vs let/const)}: \texttt{var} ha visibilità di funzione (Function Scope). \texttt{let} e \texttt{const} hanno visibilità di blocco (Block Scope, es. visibili solo dentro l'if o il for).
\end{itemize}

\begin{jsbox}[title={Snippet: Dichiarazioni e Closures}]
// Arrow Function (Brevi e `this` lessicale)
const somma = (a, b) => a + b; 

// Closure: Incapsulamento / "Memoria" locale
function counter(start) {
    let count = start; // Privato grazie al Block Scope della funzione
    return () => ++count; // Può leggere 'count' perché è nello stesso blocco lessicale
}
\end{jsbox}

\subsection{Lezioni 08-09: Stringhe, JSON, Insiemi e Math}
Le stringhe sono immutabili (metodi come \texttt{toUpperCase}, \texttt{replace(a, b)}, \texttt{split(",")} restituiscono NUOVE stringhe).

\begin{jsbox}[title={Snippet: Math, Set e JSON}]
// Oggetto Math
Math.max(1, 10, 5); // 10
Math.floor(5.9); // 5 (pavimento)
Math.random(); // Numero casuale tra [0, 1)

// Oggetto Set (Valori Univoci in memoria continua rapida)
let s = new Set([1, 2, 2, 1]); // Set {1, 2}
s.add(3); s.delete(1); s.has(2); // true

// Serializzazione in JSON
let str = JSON.stringify({a: 1}); // '{"a":1}'
let x = JSON.parse(str); // {a:1} -> NUOVO oggetto. Le FUNZIONI e DATE si perdono nel parse!
\end{jsbox}

\subsection{Extra: Espressioni Regolari (RegEx)}
Potentissime per ricerca, convalida e manipolazione di stringhe.
Sintassi letterale: \texttt{/pattern/flags} (es. flag \texttt{g} per globale, \texttt{i} per case-insensitive).

\begin{center}
\renewcommand{\arraystretch}{1.5}
\begin{tabular}{|l|l|l|}
\hline
\textbf{Classe} & \textbf{Metodo} & \textbf{Descrizione} \\
\hline
\texttt{RegExp} & \texttt{test(str)} & Vero se la stringa contiene il pattern. (Veloce!) \\
\hline
\texttt{RegExp} & \texttt{exec(str)} & Restituisce un Array con i dettagli del match, o \texttt{null}. \\
\hline
\texttt{String} & \texttt{match(rgx)} & Come exec, ma se flag \texttt{g} restituisce tutti i match. \\
\hline
\texttt{String} & \texttt{replace(rgx, n)} & Sostituisce i match con \texttt{n}. \\
\hline
\end{tabular}
\end{center}

\begin{jsbox}[title={Snippet: Utilizzo Pratico di RegEx}]
const regexEmail = /^[^\s@]+@[^\s@]+\.[^\s@]+$/; 
regexEmail.test("test@email.com"); // true

const testo = "Il numero è 123 e poi 456";
const regexNum = /\d+/g; // \d = cifre numeriche, + = una o più, g = tutte le occorrenze
console.log(testo.match(regexNum)); // ["123", "456"]

// Replace avanzato rimpiazzando tutte le vocali
console.log("ciao".replace(/[aeiou]/ig, "*")); // "c**o"
\end{jsbox}

\subsection{Extra: Metodi Utili e Universali (Da Conoscere a Memoria)}
Metodi built-in globali o di classi base che salvano spesso la vita.

\begin{center}
\renewcommand{\arraystretch}{1.5}
\begin{tabular}{|p{5.5cm}|p{8cm}|}
\hline
\textbf{Metodo} & \textbf{Comportamento} \\
\hline
\texttt{Number.parseInt(str, 10)} & Parsa una stringa in Intero (Base 10). Se fallisce: \texttt{NaN}. \\
\hline
\texttt{Number.parseFloat(str)} & Parsa stringhe in numeri a virgola mobile. \\
\hline
\texttt{isNaN(val)} o \texttt{Number.isNaN(val)} & Controlla se un valore è "Not a Number". (Il secondo è più rigoroso). \\
\hline
\texttt{Array.isArray(obj)} & Unico modo sicuro per capire se un Oggetto è in realtà un Array. \\
\hline
\texttt{Object.hasOwn(obj, "prp")} & (Moderno vs \texttt{hasOwnProperty}). Controlla se la proprietà appartiene DIRETTAMENTE all'oggetto, ignorando i prototipi. \\
\hline
\texttt{String.prototype.padStart(l, c)} & Aggiunge caratteri \texttt{c} in testa finché lunga \texttt{l} (\texttt{"5".padStart(2, "0")} $\to$ \texttt{"05"}). \\
\hline
\texttt{String.prototype.trim()} & Rimuove gli spazi/a capo vuoti a inizio e fine. \\
\hline
\end{tabular}
\end{center}

%============================
% PART 2: Lezioni 11-18
%============================
\subsection{Lezioni 10-11: Callback e Ricorsione}
\begin{itemize}
    \item \textbf{Callback}: Una funzione passata come argomento ad un'altra funzione per essere eseguita in un secondo momento. Esempio classico: \texttt{setTimeout(fn, ms)}.
    \item \textbf{Ricorsione}: Una funzione che chiama sé stessa. Regole d'oro per evitare lo \textit{Stack Overflow}:
    \begin{enumerate}
        \item \textbf{Caso Base}: Condizione di terminazione in cui si restituisce un valore senza rincorrere la funzione.
        \item \textbf{Passo Ricorsivo}: Chiamata a sé stessa ma con un "problema più piccolo" (es. \texttt{n - 1}).
    \end{enumerate}
\end{itemize}
\begin{jsbox}[title={Snippet: Ricorsione Base (Fattoriale)}]
function fattoriale(n) {
    if (n <= 1) return 1; // Caso Base
    return n * fattoriale(n - 1); // Passo Ricorsivo
}
\end{jsbox}

\subsection{Lezioni 12-14: Liste Concatenate (Linked Lists)}
Strutture dati dinamiche in cui ogni nodo punta al successivo.

\begin{jsbox}[title={Snippet: Definizione e Ricorsione su Liste}]
// Nodo Base
function NodoList(val, next = null) {
    this.val = val; 
    this.next = next;
}
let lista = new NodoList(1, new NodoList(2, new NodoList(3))); // 1 -> 2 -> 3

// Calcolare la lunghezza (Ricorsivo)
function length(l) {
    if (!l) return 0;
    return 1 + length(l.next);
}

// Stampa tutti gli elementi (Ricorsivo)
function printList(l) {
    if (!l) return;
    console.log(l.val);
    printList(l.next);
}
\end{jsbox}

\subsection{Lezioni 16-17: Array Avanzati e Funzioni di Ordine Superiore}
Metodi costosi ma potentissimi. NON mutanti (tranne \texttt{sort} o simili esplicitati), restituiscono copie calcolate e permettono composizione.

\begin{center}
\renewcommand{\arraystretch}{1.5}
\begin{tabular}{|l|p{5cm}|l|}
\hline
\textbf{Metodo (HOF)} & \textbf{Descrizione} & \textbf{Ritorna} \\
\hline
\texttt{forEach(fn)} & Esegue \texttt{fn} su ogni elemento \textit{(Sostituto del for-loop)}. & \texttt{undefined} \\
\hline
\texttt{map(fn)} & Trasforma ogni elemento. \textit{(Stessa lunghezza!)} & \textbf{Nuovo Array} \\
\hline
\texttt{filter(fn)} & Mantiene solo elementi per cui \texttt{fn===true}. & \textbf{Nuovo Array} \\
\hline
\texttt{reduce(fn, acc)} & Accumula in un singolo valore. ES: \texttt{(acc, e)=>(acc+e)}. & \textbf{Valore / Array / Obj} \\
\hline
\texttt{every(fn)} / \texttt{some(fn)} & Rispettivamente "Tutti" e "Almeno 1" passano il test \texttt{fn}. & Booleano (\texttt{true/false}) \\
\hline
\texttt{find(fn)} & Trova il \textbf{primo} elemento che passa la funzione \texttt{fn}. & Elemento o \texttt{undefined} \\
\hline
\texttt{findIndex(fn)} & Trova l'\textbf{indice} del primo elemento che passa \texttt{fn}. & Indice numerico o \texttt{-1} \\
\hline
\end{tabular}
\end{center}

\begin{jsbox}[title={Snippet: Array Chaining e Reduce Avanzato}]
let a = [1, 2, 3, 4, -5];

// Array Chaining (Concatenare i metodi restituenti Array)
let r = a.filter(v => v > 0).map(v => v * 10); // [10, 20, 30, 40]

// Somma con reduce
let somma = a.reduce((acc, curr) => acc + curr, 0); 
\end{jsbox}

\subsection{Lezione 18: Alberi Binari}
\begin{jsbox}[title={Snippet: Nodo e Visite (Pre, In, Post-Order)}]
function NodoBin(val, sx = null, dx = null) {
    this.val = val; this.sx = sx; this.dx = dx;
}

// Visita Ricorsiva Pre-Order
function dfs(root) {
    if (!root) return;
    console.log(root.val); // Root
    dfs(root.sx);          // Left
    dfs(root.dx);          // Right
}
// Visita BFS (Ampiezza / Iterativa tramite Coda)
function bfs(root) {
    let coda = [root];
    while(coda.length > 0) {
        let n = coda.shift(); // Estrae dalla testa della coda
        if(n) {
            console.log(n.val);
            if(n.sx) coda.push(n.sx); // Accoda
            if(n.dx) coda.push(n.dx);
        }
    }
}
\end{jsbox}

%============================
% PART 3: Lezioni 20-28
%============================
\subsection{Lezioni 20-21: Grafi e Reti}
Rappresentazione tramite \textbf{Liste di Adiacenza} (Mappe o Oggetti in cui la chiave è il nodo e il valore è l'array dei vicini).
\begin{jsbox}[title={Snippet: Visita Grafo (DFS) con Set dei Visitati}]
const grafo = { 'A': ['B','C'], 'B': ['A'], 'C': ['A'] };

function dfs_grafo(nodo, visitati = new Set()) {
    if (visitati.has(nodo)) return; // Ciclo evitato!
    console.log(nodo);
    visitati.add(nodo);
    
    for (let vicino of grafo[nodo]) {
        dfs_grafo(vicino, visitati);
    }
}
\end{jsbox}

\subsection{Lezioni 22-26: Identità Oggetti, Prototipi, Classi e Moduli}
\begin{itemize}
    \item \textbf{Identity}: Due oggetti sono uguali (\texttt{obj1 === obj2}) SOLO se occupano la stessa locazione di memoria, NON se hanno le stesse proprietà.
    \item \textbf{Prototipi e \texttt{new}}: Ogni oggetto ha una proprietà nascosta \texttt{\_\_proto\_\_}. Quando usiamo \texttt{new Funzione()}, JavaScript crea un oggetto e gli assegna come \texttt{\_\_proto\_\_} la proprietà \texttt{Funzione.prototype}.
    \item \textbf{Accessors System}: I \texttt{get} e \texttt{set} permettono di nascondere computazioni complesse dietro al mascheramento di semplici attributi.
    \item \textbf{Moduli}: Divisione in file (\texttt{export} e \texttt{import}).
\end{itemize}

\begin{jsbox}[title={Snippet: Classi (ES6) ed Ereditarietà}]
class Animale {
    constructor(nome) { this.nome = nome; }
    verso() { console.log("..."); }
}
class Cane extends Animale {
    static specie = "Canis lupus familiaris"; // Statico, appartiene a 'Cane'
    #fame = 100;                              // Privato, accessibile solo qui dentro
    
    constructor(nome, razza) {
        super(nome); // Obbligatorio! Richiama parent constructor
        this.razza = razza;
    }
    
    get fame() { return this.#fame; } // Getter
    verso() { console.log("Bau"); } // Overriding del metodo
}
\end{jsbox}

\subsection{Lezione 27: Gestione delle Eccezioni}
Permette di riprendersi da errori letali o segnalare inconsistenze (es. parametri sbagliati in ingresso).
\begin{jsbox}[title={Snippet: Costrutto Try-Catch-Finally ed Errori Custom}]
class ParametroErrato extends Error { } // Custom Error

function calcolaArea(lato) {
    if (lato <= 0) throw new ParametroErrato("Il lato deve essere positivo"); // Lancia eccezione
    return lato * lato;
}

try {
    let r = calcolaArea(-5);
} catch (e) {
    if (e instanceof ParametroErrato) console.log("Errore specifico: " + e.message);
    else throw e; // Rilancia errori inaspettati (es. Syntax Error)
} finally {
    console.log("Eseguito sempre, errore o meno. Ottimo per chiudere i file.");
}
\end{jsbox}

\subsection{Lezione 28: Generatori ed Utilities Esercitazioni}
I generatori sono funzioni che possono essere messe in pausa (\texttt{yield}) e riprese (\texttt{next()}).

\begin{jsbox}[title={Snippet: Function* e Yield}]
function* range(min, max) {
    for (let i = min; i <= max; i++) {
        yield i; // Pausa e restituisci 'i'
    }
}
// Utilizzo (Iteratore lazy):
for (let n of range(1, 3)) console.log(n); // Stampa 1, 2, 3
\end{jsbox}

\ifx\isStandalone\undefined
\end{document}
\fi
